\chapter{حجية خبر الآحاد والأدلة على قبوله}

\section{نضارة وجه مبلغ السنة ودلالتها على الآحاد}
الحمد لله والصلاة والسلام على رسول الله، أما بعد؛ فقد بوب الإمام الشافعي رحمه الله تعالى باباً عظيماً لبيان الحجة في تثبيت خبر الواحد، ليرد به على المعتزلة وأهل البدع الذين يشككون في خبر الواحد ويردونه في العقائد أو الأحكام.
استدل الشافعي أولاً بحديث عبد الله بن مسعود رضي الله عنه أن النبي ﷺ قال: «نَضَّرَ الله عبداً -وفي رواية: امرأً- سَمِعَ مَقَالَتي فَحَفِظَهَا وَوَعَاهَا وَأَدَّاهَا، فَرُبَّ حَامِلِ فِقْهٍ غَيْرُ فَقِيهٍ، وَرُبَّ حَامِلِ فِقْهٍ إِلَى مَنْ هُوَ أَفْقَهُ مِنْهُ. ثَلَاثٌ لا يُغِلُّ عَلَيْهِنَّ قَلْبُ مُسْلِمٍ: إِخْلَاصُ العَمَلِ للهِ، وَالنَّصِيحَةُ لِلْمُسْلِمِينَ، وَلُزُومُ جَمَاعَتِهِمْ، فَإِنَّ دَعْوَتَهُمْ تُحِيطُ مِنْ وَرَائِهِمْ».
والنضارة هي حسن الوجه وبريقه. وقد استنبط الشافعي من هذا الحديث حجية خبر الواحد؛ لأن النبي ﷺ ندب إلى استماع مقالته وحفظها وأدائها، وخص بالذكر «امرأً» أو «عبداً» وهو (الواحد). فدل ذلك على أنه لا يأمر الواحد بأن يؤدي عنه إلا وخبر هذا الواحد تقوم به الحجة على من أُدي إليه؛ لأنه إنما يؤدي حلالاً وحراماً وأمراً ونهياً. 
كما دل الحديث على التفرقة بين الحفظ والفقه، فليس كل راوٍ وحافظ فقيهاً مستنبطاً، والخير كل الخير فيمن جمع بين الحفظ والفقه والأصول.

\section{لزوم جماعة المسلمين وحجية الإجماع}
في قوله ﷺ: «ولزم جماعتهم» حجة ظاهرة على اعتبار إجماع المسلمين حجة لازمة وملزمة، لا يجوز لمؤمن يؤمن بالله واليوم الآخر أن يخالفه. وما أجمع عليه الصحابة والتابعون سلفاً لا يجوز نقضه بادعاء إجماعات معاصرة مبنية على شريعة مستوردة كالديمقراطية أو غيرها، فمن ادعى الإجماع على باطل يخالف إجماع السلف الصالح المتقدم فهو كمن ادعى إجماع الرافضة على تكفير الصحابة؛ كلاهما إجماع باطل لا وزن له في دين الله.

\section{وجوب التسليم للسنة (حديث الأريكة)}
ساق الشافعي حديث أبي رافع رضي الله عنه أن النبي ﷺ قال: «لَا أُلْفِيَنَّ أَحَدَكُمْ مُتَّكِئًا عَلَى أَرِيكَتِهِ يَأْتِيهِ الأَمْرُ مِنْ أَمْرِي مِمَّا أَمَرْتُ بِهِ أَوْ نَهَيْتُ عَنْهُ، فَيَقُولُ: لَا نَدْرِي! مَا وَجَدْنَا فِي كِتَابِ اللهِ اتَّبَعْنَاهُ». 
وهذا تثبيت من النبي ﷺ لسنته وإعلام للأمة أنها لازمة لهم وإن لم يجدوا لها نص حكم في كتاب الله. فالأريكة هي السرير أو ما يتكئ عليه الرجل، وفيه وعيد شديد لمن يرد سنة رسول الله ﷺ بحجة الاكتفاء بالقرآن، فمن رد السنة فقد رد أمر الله جل وعلا.

\section{أمثلة وتطبيقات من العهد النبوي لقبول خبر الواحد}
ذكر الشافعي أمثلة عملية متواترة المعنى تدل على أن الصحابة كانوا يقبلون خبر الواحد ويقيمون به الحجة وينقضون به اليقين السابق:

\subsection{خبر أم سلمة في القبلة للصائم}
رجل من الأنصار قَبَّل امرأته وهو صائم، فوجد في نفسه من ذلك، فأرسل امرأته تسأل أم سلمة رضي الله عنها، فأخبرتها أن النبي ﷺ يقبل وهو صائم. فرجعت إلى زوجها فلم يقبل ظناً منه أنها خصوصية للنبي ﷺ. فلما استوثق وعلم، سَلَّم. والشاهد أن النبي ﷺ قال لأم سلمة: «أَلَا أَخْبَرْتِيهَا أَنِّي أَفْعَلُ ذَلِكَ؟» فلو لم تكن الحجة قائمة بخبر امرأة واحدة عن النبي ﷺ لم يأمرها بذلك.

\subsection{تحول أهل قباء عن قبلتهم بخبر واحد}
كان الصحابة يصلون الصبح في قباء متوجهين إلى بيت المقدس (وهي القبلة اليقينية التي فرضوها وسمعوها من النبي ﷺ)، فجاءهم آتٍ واحد، فأخبرهم أن القبلة قد حولت إلى الكعبة، فاستداروا وهم في صلاتهم. فلو لم يكن خبر الواحد الثقة حجة قاطعة، ما ترك الصحابة فرضاً متيقناً لخبر رجل واحد، وما استداروا دون أن يرسلوا ليتأكدوا من النبي ﷺ. فلما أقرهم النبي ﷺ على ذلك، ثبت أن خبر الواحد حجة.

\subsection{إراقة أبي طلحة للخمر بخبر المخبر}
كان أنس بن مالك رضي الله عنه يسقي أبا طلحة وأبا عبيدة بن الجراح وأبي بن كعب شراباً من فضيخ وتمر (وكان حلالاً حينها). فجاءهم رجل واحد فقال: إن الخمر قد حُرِّمت. فقال أبو طلحة: يا أنس، قم إلى هذه الجرار فاكسرها. 
ومعلوم أن إضاعة المال محرمة، وهؤلاء من كبار الصحابة وعلماء الأمة، فلو لم تقم الحجة عندهم بخبر هذا الواحد ما أراقوا هذا المال الذي كان حلالاً إلى وقت قريب، وما رضوا بإتلافه بمجرد خبر رجل واحد.

\subsection{إرسال الآحاد لتبليغ الأحكام والحدود والعقائد}
\begin{itemize}
    \item \textbf{في الحدود:} أرسل النبي ﷺ أنيساً الأسلمي وحده في قضية العسيف قائلاً: «وَاغْدُ يَا أُنَيْسُ إِلَى امْرَأَةِ هَذَا، فَإِنْ اعْتَرَفَتْ فَارْجُمْهَا». فلو لم تقم الحجة بالواحد ما أرسله وحده لإقامة حد الرجم الذي يزهق النفس.
    \item \textbf{في الأحكام ومناسك الحج:} أرسل علي بن أبي طالب رضي الله عنه وحده على جمل يصرخ في منى: «إنها أيام طعام وشراب فلا يصومن أحد». وأرسل ابن مربع الأنصاري وحده ليقول للناس في عرفة: «قِفُوا عَلَى مَشَاعِرِكُمْ فَإِنَّكُمْ عَلَى إِرْثٍ مِنْ إِرْثِ أَبِيكُمْ إِبْرَاهِيمَ». وبعث أبا بكر وعلياً رضي الله عنهما وحدهما لإقامة الحج ونبذ العهود للمشركين.
    \item \textbf{في العقائد والتوحيد:} بعث معاذ بن جبل وحده إلى اليمن ليدعوهم إلى التوحيد (شهادة أن لا إله إلا الله)، وإلى الصلاة والزكاة. فلو كان خبر الواحد لا يُقبل في العقائد (كما تزعم المبتدعة) لما أرسل النبي ﷺ معاذًا وحده لدعوة أهل اليمن للعقيدة والأحكام معاً.
\end{itemize}

\section{الفرق بين القاضي والشاهد في الأخبار}
أوضح الشافعي أن قضاء القاضي هو في حقيقته إخبار وإمضاء لبينة أو إقرار ثبت عنده. فالقاضي يلزمه أن ينفذ حكم الله بعلمه من خلال البينة، ويكون مخبراً بالحلال والحرام الذي لزمه. بينما الشاهد الذي يشهد على واقعة لا ينفذ الحكم بشهادته منفرداً، بل لا بد من نصاب الشهادة (كشاهدين). فدل هذا على أن مقام النقل والرواية (كالإخبار بالحكم الشرعي) يختلف عن مقام الشهادة على الوقائع والأعيان.